
%% =========================================
%%   DEL 1 - UTEN HJELPEMIDLER
%% =========================================

\Del{u}

\Oppgave[2] % Oppgave 1.1
\points*{2}

I en kasse ligger det $60$ epler. $\SI{20}{\percent}$ av eplene er grønne,
$\cfrac{7}{12}$ er røde og resten er gule. Hvor mange epler er gule?


\Oppgave[2] % Oppgave 1.2
\points*{2}

Regn ut og skriv svaret på standardform
%
\begin{equation*}
    \frac{\num{7.5E5} \cdot \num{4.0E6}}{\num{2.5E-4}}
\end{equation*}


\Oppgave[2] % Oppgave 1.3
\points*{2}

Svein hadde en nominell lønn på $\num{440 000}$ i $2008$. Konsumprisindeksen var
$88$ dette året.
\bigskip\noindent

Bestem reallønna til Svein i $2008$.


\Oppgave[3] % Oppgave 1.4
\points*{3}

Ved en skole er det $150$ ansatte. Skolen har et treningsrom som alle ansatte
kan benytte. Tabellen nedenfor viser hvor ofte de ansatte benyttet
treningsrommet i løpet av en uke.

\begin{table}[H]
    \centering
    \caption{}
    \begin{tabular}{| S[table-format=2.0] | S[table-format=2.0] |}
         \tableHeaders{Antall ganger}{Antall ansatte}
         0 & 80 \\
         1 & 10 \\
         2 & 15 \\
         3 & 25 \\
         4 &  5 \\
         5 & 15 \\ \hline
    \end{tabular}
    \label{tab:del-1-oppgave-1.4}
\end{table}

\Oppgave[2] % Oppgave 1.5
\points*{2}

\begin{figure}[H]
    \centering
    \tikzsetnextfilename{Forkurs-1p-2p-laererutdanning-2018-V-oppgave-1-5}
    \begin{tikzpicture}
        \edef\BCxAB{2}
        \edef\ArealABC{10}
        \pgfmathsetmacro{\AB}{sqrt(2*\ArealABC/\BCxAB)}
        \pgfmathsetmacro{\BC}{\BCxAB*\AB}
        \tkzDefPoint(0,0){A}
        \tkzDefPoint(\AB,0){B}
        \tkzDefPoint(\AB,\BC){C}

        \tkzMarkRightAngle[scale=2](A,B,C)
        \tkzDrawPolygon(A,B,C)

        \tkzLabelPoint[below left](A){$A$}
        \tkzLabelPoint[below right](B){$B$}
        \tkzLabelPoint[above right](C){$C$}
    \end{tikzpicture}
    \caption{}
    \label{fig:del-1-oppgave-1-5}
\end{figure}

Arealet av en trekant er $\SI{49}{\cm\squared}$, $BC$ er dobbelt så langt som
$AB$. Se skissen ovenfor. Bestem lengden av $AB$.

\Oppgave[2] % Oppgave 1.6

\begin{figure}[H]
    \centering
    \tikzsetnextfilename{Forkurs-1p-2p-laererutdanning-2018-V-oppgave-1-6}
    \begin{tikzpicture}
        \begin{axis}[
            Eksamen1,
            ytick={0,1,...,8},
            yticklabel = {\num{\fpeval{\tick*1000}}},
            xtick={0,2,...,20},
            ymin=-.1,
            ymax=9,
            xmin=-.1,
            xmax=22,
            domain = 0.5:25,
        ]
            \addplot[color=maincolorMedium,thick,samples=250] {8/x} node[right,pos=0.25] {$f$};
            \node[] at (axis cs: 6,8.5) {Pris per elev (kroner)};
            \node[] at (axis cs: 19,1.5) {Antall elever};

            \draw [maincolorMedium, fill] (axis cs: 2,4) circle (3pt) node [above right] {$(2, \num{4000})$};
            \draw [maincolorMedium, fill] (axis cs: 4,2) circle (3pt) node [above right] {$(4, \num{2000})$};
            \draw [maincolorMedium, fill] (axis cs: 10,0.8) circle (3pt) node [above right] {$(10, \num{800})$};
            \draw [maincolorMedium, fill] (axis cs: 16,0.5) circle (3pt) node [above right] {$(16, \num{500})$};
        \end{axis}
    \end{tikzpicture}
    \caption{}
    \label{fig:del-1-oppgave-1-6}
\end{figure}

Elevene i klasse $2$A vil leie et lokale for å arrangere klassefest. De vil
spleise på utgiftene. Ovenfor ser du grafen til en funksjon $f$. Grafen viser
sammenhengen mellom hvor mange elever som blir med på festen, og prise hver elev
må betale.

\begin{oppgaver}
    \Item{1} Hvor mye må hver elev betale dersom $25$ elever blir med på festen?
\end{oppgaver}

\begin{oppgaver}
    \Item{1} Bestem funksjonsuttrykket $f(x)$.
\end{oppgaver}


\Oppgave[7] % Oppgave 1.7

I dag veier Ole $\SI{100}{\kg}$. Han har som mål å gå ned $\SI{16}{\kg}$ i vekt
i løpet av de neste $40$ ukene. \bigskip

Anta at vektnedgangen er lineær, og at Ole akkurat når målet sitt.

\begin{oppgaver}
    \Item{2} Sett opp en modell som viser hvor mye Ole vil veie om $x$ uker
    dersom antakelsene er riktig.
    \label{delopg:del-7-oppgave-1a}
\end{oppgaver}

\begin{oppgaver}
    \Item{1} Hvor mange uker vil det gå før Ole veier $\SI{88}{\kg}$, i følge
    modellen fra \cref{delopg:del-7-oppgave-1a}?
\end{oppgaver}

Ole vil gjerne lage en eksponentiell modell for vektnedgangen. Han antar vekten
vil gå ned med $\SI{0.4}{\percent}$ hver uke.

\begin{oppgaver}
    \Item{2} Sett opp en eksponentiell modell ut i fra Ole's antakelse.
    \label{delopg:del-7-oppgave-1c}
\end{oppgaver}

\begin{oppgaver}
    \Item{2} I hvilken uke vil Ole gå mest ned i vekt, ifølge modellen fra
    \cref{delopg:del-7-oppgave-1c}?
    Vil han nå målet sitt hvis vektnedgangen følger denne modellen?
\end{oppgaver}


\Oppgave[4] % Oppgave 1.8

\begin{figure}[H]
    \centering
    \begin{subfigure}[b]{0.32\textwidth}
        \centering
        \tikzsetnextfilename{Forkurs-1p-2p-laererutdanning-2018-V-oppgave-1-8a}
        \polygons{1}
        \caption{}
        \label{fig:del-1-oppgave-8-a}
    \end{subfigure}\hfill%
    \begin{subfigure}[b]{0.32\textwidth}
        \centering
        \tikzsetnextfilename{Forkurs-1p-2p-laererutdanning-2018-V-oppgave-1-8b}
        \polygons{2}
        \caption{}
        \label{fig:del-1-oppgave-8-b}
    \end{subfigure}\hfill%
    \begin{subfigure}[b]{0.32\textwidth}
        \centering
        \tikzsetnextfilename{Forkurs-1p-2p-laererutdanning-2018-V-oppgave-1-8c}
        \polygons{3}\vspace{-0.45cm}
        \caption{}
        \label{fig:del-1-oppgave-8-c}
    \end{subfigure}
    \caption{}\label{fig:del-1-oppgave-8}
\end{figure}

I \cref{fig:del-1-oppgave-8} ser du tre figurer. Figurene er satt sammen av små
sirkler. Tenk deg at du skal fortsette å lage figurer etter samme mønster.

\begin{oppgaver}
    \Item{3} Skriv av tabellen nedenfor, og fyll ut det som mangler. Gjør
    beregninger, eller forklar hvordan du tenker.
\end{oppgaver}

\begin{table}[H]
    \centering
    \caption{}
    \begin{tabular}{| S[table-format=1.0] | S[table-format=2.0] |}
      \tableHeaders{Figur}{Antall sirkler}
           1   &  6 \\
           2   & 10 \\
           3   & 14 \\
           4   &    \\
           5   &    \\
         {$n$} &    \\ \hline
    \end{tabular}
    \label{tab:del-1-oppgave-1.8}
\end{table}

\begin{oppgaver}
    \Item{1} Hvor mange sirkler vil det være i figur nummer $100$?
\end{oppgaver}

